\documentclass[10pt,a4paper]{amsart}
\usepackage[margin=19mm]{geometry}
\usepackage{amsmath,amssymb,mathtools}
\usepackage[scr=rsfs,cal=euler]{mathalfa}
\usepackage{xfrac}
\usepackage[unicode,hidelinks,bookmarks=false]{hyperref}
\newcommand{\ie}{i.e.\ }
\newcommand{\cf}{cf.\ }
\newcommand{\resp}{respectively\ }
\newcommand{\Subsection}{\subsection}
\providecommand{\nobreakdash}{\nobreak\hbox{-}}
\setlength{\emergencystretch}{2em}
\multlinegap0pt
\usepackage{amssymb}
\newtheorem{lemma}{Lemma}
\title{Almost Every Graph Is Reconstructible from Its Token Graphs}
\author{Yuanming Luo}
\date{Source: arXiv:2608.22094v2 (CC BY 4.0)}
\begin{document}
\maketitle

\noindent\emph{Source and licence.} Almost Every Graph Is Reconstructible from Its Token Graphs. Version arXiv:2608.22094v2 (CC BY 4.0), distributed by arXiv under the Creative Commons Attribution 4.0 International licence, http://creativecommons.org/licenses/by/4.0/, as recorded in the arXiv licence metadata for this exact version and shown on the abstract page https://arxiv.org/abs/2608.22094v2. Full licence notice: \texttt{licenses/arxiv-2608.22094-CC-BY-4.0.txt}.\par\smallskip

\noindent\textit{Editorial scope.} Counted unit: the article's lemma lem:commutation (source0.tex lines 36--38). The article's complete original proof of it is delivered with it (source0.tex lines 40--48, one \texttt{proof} environment of 9 lines). No mathematics is written, repaired or re-derived: the statement and the proof are the article's own lines, in the article's own notation, and no retained line is substituted or re-flowed.  No further source-local context is needed: every cross-reference of every retained line resolves inside the two delivered spans.  Nothing was omitted from any retained span, so the package carries no omission record.  No retained line contains a citation, so the wrapper carries no bibliography and no bibliography entry.  No retained line contains a figure environment, an image-inclusion command or drawing code, so no image is delivered and the package declares no asset dependency.  Editorial formatting only, in addition: an AMS \texttt{amsart} preamble that restates the article's own declarations and macros used by the retained lines, namely the environments \texttt{lemma} on separate counters, together with the standard packages those lines need.  The retained environments print in this document as Lemma 1 in order of appearance, whereas the article prints the same items with the numbering of its own full document.  The complete native source of the article as distributed in the arXiv e-print for this exact version is bundled unchanged as \texttt{sources/208302/source0.tex}.
\begin{lemma}\label{lem:commutation}
For every graph $G$ and every $1\leq k<|V(G)|$, we have $C(F_k(G))=F_k(C(G))$.
\end{lemma}

\begin{proof}
Let $U,V\in V(F_k(G))$ be distinct. Suppose $U\triangle V=\{u,v\}$. Write $R=U\cap V$, so $U=R\cup\{u\}$ and $V=R\cup\{v\}$. Let $w$ be a common neighbor of $u$ and $v$ in $G$. If $w\notin R$, then $R\cup\{w\}$ is adjacent to both $U$ and $V$ in $F_k(G)$. If $w\in R$, then $(R\setminus\{w\})\cup\{u,v\}$ is adjacent to both $U$ and $V$. Thus every common neighbor of $u$ and $v$ gives a common neighbor of $U$ and $V$.

Conversely, let $W$ be a common neighbor of $U$ and $V$ in $F_k(G)$. Since $W$ differs from both $U$ and $V$ by one token move, either $W=R\cup\{w\}$ for some $w\notin R\cup\{u,v\}$, or $W=(R\setminus\{w\})\cup\{u,v\}$ for some $w\in R$. In the first case, adjacency of $W$ to $U$ and $V$ implies $uw,vw\in E(G)$. The second case gives the same conclusion. Hence $w$ is a common neighbor of $u$ and $v$.

Thus $U$ and $V$ have the same number of common neighbors in $F_k(G)$ as $u$ and $v$ have in $G$. Moreover, $U$ and $V$ are adjacent in $F_k(G)$ if and only if $u$ and $v$ are adjacent in $G$. Therefore $\{U,V\}\in E(C(F_k(G)))$ if and only if $\{u,v\}\in E(C(G))$.

Let $n=|V(G)|$. Now suppose $|U\triangle V|\geq 4$. If $|U\triangle V|=4$, then $U$ and $V$ have at most four common neighbors in the Johnson graph $J(n,k)$, and therefore at most four common neighbors in $F_k(G)$. If $|U\triangle V|\geq 6$, they have no common neighbor in $J(n,k)$, since two token moves can relocate at most two tokens. Thus $C(F_k(G))$ does not add an edge between $U$ and $V$. They are also nonadjacent in $F_k(C(G))$ because their symmetric difference has size at least four. Hence the two graphs have the same edge set.
\end{proof}

\end{document}
